\mode*
\section{Search and verification protocol}
\label{app:protocol}

This appendix documents the literature searches behind the references, for
reproducibility, following the protocol of the companion paper
\autocite[App.~literature protocol]{VTProgMisconceptions}: every reference
was found by a recorded query, read at least at abstract level, and
recorded with a provenance block (claim, query, selection rationale,
verbatim supporting quote, verification level) directly above its entry in
\texttt{misconceptions.bib}.

\subsection{Phase one: seed round (2026-07-18)}

The seed round was a claim-driven search conducted with the
\texttt{scholar} tool against DBLP, OpenAlex, and Semantic Scholar
(Semantic Scholar was rate-limited during part of the session), with
abstracts fetched via the OpenAlex API.
The queries whose results were retained for this paper:

\begin{itemize}
  \item \texttt{scholar search "latex word document preparation comparison" -p s2 -p openalex}
  \item \texttt{scholar search "black box inside glass box presenting computing concepts" -p openalex -p dblp}
  \item \texttt{scholar search "Sorva notional machines introductory programming" -p dblp -p openalex}
  \item \texttt{scholar search "Ben-Ari constructivism" -p dblp}
  \item \texttt{scholar search "difficulties of learning to program du Boulay" -p openalex}
  \item \texttt{scholar search '"variation theory" computer programming' -p openalex}
\end{itemize}

The central negative result of this round: no studies of LaTeX-specific
learner misconceptions were found; the LaTeX-related hits beyond
\autocite{Knauff2014Latex} were off-topic (accessibility, general
tutorials, or chemistry).
Not yet verified at full text (tracked in the course repository's issues
\#124--\#125): \autocite{DuBoulay1981BlackBox},
\autocite{BenAri1998Constructivism}.

\subsection{Phase two: systematic round}
\label{sec:phase-two}

The systematic round (2026-07-18) was run with the \texttt{scholar} tool
against DBLP, OpenAlex and Semantic Scholar (\texttt{s2}), with abstracts
read from OpenAlex (via \texttt{scholar} and, for one source, the Semantic
Scholar Graph API).
Because phase one established that the LaTeX-specific learner literature is
empty, this round targeted five adjacent or repeat areas:
(a)~markup languages in education;
(b)~comprehension of compiler and programming error messages;
(c)~end-user programming of documents;
(d)~WYSIWYG-versus-source editing; and
(e)~a repeat attempt at LaTeX-specific learner studies.
Screening kept a source only if its abstract, read in full, \emph{entailed}
a claim used in the text at the strength the source supports; adjacent-%
literature sources were additionally required to yield a finding
transferable to LaTeX, and the chapter prose marks every such transfer as a
conjecture.

The queries that retained sources:
\begin{itemize}
  \item \texttt{scholar search "compiler error messages novice programmers
    comprehension" -p dblp -p openalex} (b)---retained
    \autocite{Munson2016CompilerErrors,Traver2010CompilerErrors,%
    Prather2017Interaction}, and \autocite{Ko2004Barriers} (c), which also
    surfaced here.
  \item \texttt{scholar search "enhanced error messages students
    introductory programming" -p dblp -p openalex} (b)---retained
    \autocite{Becker2016CEM,Denny2021ErrorReadability}.
  \item \texttt{scholar search "teaching HTML students misconceptions web
    development" -p s2} (a)---retained \autocite{Kharchenko2025WebDesign}.
\end{itemize}

The queries that returned nothing retainable, recorded for reproducibility:
\begin{itemize}
  \item Repeat LaTeX-specific attempts (e):
    \texttt{scholar search "learning LaTeX typesetting students novice" -p
    dblp -p openalex} and \texttt{scholar search "learning teaching LaTeX
    students difficulties" -p dblp -p openalex} returned only
    \autocite{Knauff2014Latex} (already held) and a non-empirical
    LaTeX-workshop experience report; no learner-difficulty study.
  \item Markup in education (a):
    \texttt{scholar search "teaching markup language HTML education
    learning" -p dblp -p openalex},
    \texttt{"students learning HTML markup separation content
    presentation" -p openalex},
    \texttt{"teaching HTML CSS students structure presentation" -p dblp -p
    openalex} and \texttt{"markup language learning students education" -p
    openalex}---off-topic hits or empty result sets.
  \item WYSIWYG-versus-source (d):
    \texttt{scholar search "WYSIWYG versus markup document editing
    comparison" -p dblp -p openalex} and \texttt{"word processor users
    formatting styles structure novice" -p openalex}---returned tool papers
    (wikis, Markdown pipelines, HTML-article formats) but no learner study.
  \item End-user document authoring (c):
    \texttt{scholar search "end-user structured document authoring word
    processing styles" -p dblp -p openalex} and
    \texttt{"novice compilation behaviour error quotient Jadud" -p dblp -p
    openalex}---empty.
  \item Research-question-driven queries via \texttt{scholar rq} (for (b)
    and (a), e.g.\ the generated query \texttt{novice programmers learning
    difficulties understanding compiler error messages OR interpreter error
    messages}) produced no additional retainable papers: OpenAlex was
    rate-limiting (HTTP~429) during the session and DBLP returned nothing to
    classify.
\end{itemize}

Two conclusions follow.
First, the LaTeX-specific-learner-study gap is \emph{confirmed}: a second,
systematic round again found no study of how learners (mis)understand
LaTeX's compile model, errors, or semantic markup.
The paper must therefore build its catalogue from the adjacent literatures
retained above and from our own course data
(\cref{sec:method-course-data}).
Second, the adjacent literatures divide unevenly: comprehension of
programming error messages is well developed and transfers cleanly to
\cref{sec:errors}, whereas markup-in-education and end-user document
authoring are themselves thin---even for HTML and CSS
\autocite{Kharchenko2025WebDesign}---so \cref{sec:markup} leans
correspondingly harder on the course data.
